\documentclass[11pt]{article}

\usepackage[margin=1in]{geometry}
\usepackage{amsmath,amssymb}
\usepackage{booktabs}
\usepackage{threeparttable}
\usepackage[font=small,labelfont=bf]{caption}
\usepackage{enumitem}
\usepackage[hidelinks]{hyperref}

\setlist{nosep,leftmargin=*}

\title{Dantzig-Selector IV:\\Revisiting \emph{Immigration, Crime, and Crime (Mis)Perceptions}}
\author{Alex Zhao, Alexandra Zhou, Oscar Barnes, Yuao Zhou}
\date{\today}

\begin{document}
\maketitle

% =====================================================================================
\section{The paper in brief}

Ajzenman, Dominguez and Undurraga (2023, \emph{AEJ: Applied}) ask whether immigration
raises crime, and whether it raises \emph{fear} of crime. Chile is a useful setting: its
foreign-born population more than doubled in the 2010s, driven largely by arrivals from
Peru, Haiti, and Venezuela. Crime is measured with Chile's national victimization survey
(ENUSC), which records self-reported victimization and crime perceptions and so avoids
relying on police reports. Immigration is measured with residence permits by
municipality.

The main design compares changes between 2008 and 2017 across $N = 101$ municipalities.
The regressor of interest is the change in the log immigrant-to-population ratio, so each
coefficient is the effect of \emph{doubling} a municipality's immigrant share, in
percentage points. Immigrants choose where to live, so OLS may be biased. The authors
instrument with a shift-share (Bartik) instrument:
\[
  Z^{B}_m = \sum_{k=1}^{11} \theta^{k}_{m,2008} \times \Delta \ln \mathrm{MIGR}^{k},
\]
where $\theta^{k}_{m,2008}$ is nationality $k$'s share of municipality $m$'s immigrants in
2008, and $\Delta \ln \mathrm{MIGR}^{k}$ is the growth of nationality $k$'s emigrant stock
in destination countries \emph{other than Chile} (UN bilateral data, 11 countries). The
intuition is that new migrants follow earlier migrants from the same country (the
shares), and that origin-country push factors (the shifts) are unrelated to local
conditions in Chile. The headline finding is that immigration has no detectable effect on
victimization but significantly raises crime-related concerns and preventive behavior: a
widening ``crime-perception gap.''

% =====================================================================================
\section{What we did}

Goldsmith-Pinkham, Sorkin and Swift (2020) show that a Bartik estimator is equivalent to
using the 11 shares as instruments with weights fixed by the shifts, so identification
rests on the shares. We therefore drop the fixed weights, treat the 11 country shares as
\emph{separate} candidate instruments, and let the Dantzig selector choose among them in
the first stage. Let $\tilde{D}_m$ and $\tilde{\theta}_m$ denote the endogenous variable and
the shares after partialling out the controls (intercept, average age, share male), with
each share standardized. The first stage solves
\[
  \hat{\pi} = \arg\min_{\pi \in \mathbb{R}^{11}} \|\pi\|_1
  \quad \text{s.t.} \quad
  \max_{k}\left|\frac{1}{N}\sum_{m=1}^{N} \tilde{\theta}_{mk}\left(\tilde{D}_m - \tilde{\theta}_m'\pi\right)\right| \le \lambda .
\]
Partialling out the controls is equivalent to leaving them unpenalized. We set
\[
  \lambda = \hat{\sigma}\sqrt{\frac{2\log p}{N}} = 0.3445 \times \sqrt{\frac{2 \log 11}{101}} = 0.0751,
\]
the version of Candès and Tao's $\sigma\sqrt{2\log p}$ (stated for unit-norm columns) that
applies to this averaged constraint with standardized instruments. Here
$\hat{\sigma}$ is the residual standard deviation of an unpenalized OLS first stage on all
11 shares, which is feasible because $p = 11$ is small relative to $N$. We do not tune
$\lambda$ by cross-validation at this stage.

Each table below reports five estimators:
\begin{enumerate}
  \item \textbf{OLS}, as in the paper.
  \item \textbf{2SLS Bartik}: the paper's estimator, with one combined instrument.
  \item \textbf{2SLS all shares}: all 11 shares as instruments, with no selection.
  \item \textbf{Post-DS 2SLS}: 2SLS using only the shares the Dantzig selector keeps.
  \item \textbf{DS-fit IV}: just-identified IV using the Dantzig first-stage fitted value
        $\tilde{\theta}_m'\hat{\pi}$ as the single instrument.
\end{enumerate}

\paragraph{Validation.} Our code reproduces the paper's design exactly. The shares times
the shifts rebuild the paper's instrument (maximum absolute difference
$2.9 \times 10^{-8}$). The Bartik first stage gives the published $F = 17.35$ and partial
$R^2 = 0.104$, and the OLS and Bartik 2SLS point estimates in Tables~\ref{tab:concerns}
to~\ref{tab:victim} match the published Tables VI to IX.

% =====================================================================================
\section{First stage: which instruments matter?}

\begin{table}[htbp]
\centering
\small
\begin{threeparttable}
\caption{First stage: Dantzig selection among the 11 country shares}
\label{tab:first}
\begin{tabular}{lcc}
\toprule
\multicolumn{3}{l}{\emph{Panel A: Dantzig coefficients on standardized 2008 shares}}\\
\midrule
Country & Dantzig coefficient & Rotemberg weight in paper's Bartik \\
\midrule
Argentina & $-0.1214$ & --- \\
Bolivia   & 0         & 0.484 \\
Brazil    & 0         & 0.009 \\
China     & $\phantom{-}0.0084$ & 0.087 \\
Colombia  & 0         & --- \\
Ecuador   & 0         & 0.092 \\
Haiti     & 0         & --- \\
Peru      & 0         & 2.225 \\
Spain     & 0         & --- \\
USA       & 0         & --- \\
Venezuela & $-0.0209$ & --- \\
\midrule
\multicolumn{3}{l}{\emph{Panel B: Strength of each instrument set}}\\
\midrule
Instrument set & Robust first-stage $F$ & Partial $R^2$ \\
\midrule
Bartik (1 instrument, paper)        & 17.35 & 0.104 \\
All 11 shares                        & \phantom{1}7.48 & 0.373 \\
Dantzig-selected (3 shares)          & 13.81 & 0.321 \\
\bottomrule
\end{tabular}
\begin{tablenotes}[flushleft]
\footnotesize
\item Notes: Panel A reports the Dantzig first-stage coefficients, in units of the
endogenous variable (log change in the immigrant-to-population ratio) per standard
deviation of each share, after partialling out the controls. $\lambda = 0.0751$. Rotemberg
weights are from the paper, which reports only the top five; --- means not reported. Panel
B reports heteroskedasticity-robust (HC1) Wald $F$-statistics for the excluded
instruments. $N = 101$.
\end{tablenotes}
\end{threeparttable}
\end{table}

Two findings stand out in Table~\ref{tab:first}.

\emph{The fixed Bartik weights discard most of the predictive information in the shares.}
Used separately, the shares explain 37\% of the residual variation in immigration growth,
against about 10\% for the Bartik combination. The Dantzig selector keeps most of this
(32\%) with only three instruments, so its $F$-statistic (13.81) is close to the Bartik's.
Using all 11 shares without selection gives a weak first stage ($F = 7.48$). That raises
the textbook concern that many weak instruments bias 2SLS toward OLS, which is the main
reason to select.

\emph{The selected instruments are not the ones driving the paper's estimate.} The paper's
Bartik estimate relies overwhelmingly on the Peruvian share (Rotemberg weight 2.2), and its
pre-trend checks focus on Peru. The Dantzig selector instead keeps Argentina, Venezuela,
and China, and drops Peru. Argentina has the largest coefficient, and it is negative:
municipalities where Argentines made up more of the 2008 immigrant population saw
\emph{slower} immigrant growth afterward. Because shares are compositional, one reading is
that a high Argentine share stands in for a low share of the fast-growing groups. We would
not over-interpret the exact set chosen, however. When instruments are correlated, $\ell_1$
methods pick among them somewhat arbitrarily, which motivates trying the elastic net next.
The more important implication is that our estimates are identified by different
variation than the paper's, and the paper's pre-trend evidence does not directly cover it.

% =====================================================================================
\section{Second-stage results}

All tables share the same layout: the coefficient, the robust standard error in
parentheses, and the $p$-value in brackets. ``Mean'' is the outcome's sample mean across
all years, in percent.

% ---------------------------------------------------------------- Table VI
\begin{table}[htbp]
\centering
\small
\begin{threeparttable}
\caption{Crime-related concerns (paper's Table VI)}
\label{tab:concerns}
\begin{tabular}{lcccccc}
\toprule
 & & \multicolumn{2}{c}{Paper's design} & \multicolumn{3}{c}{Shares as separate instruments}\\
\cmidrule(lr){3-4}\cmidrule(lr){5-7}
Outcome & Mean & OLS & 2SLS Bartik & 2SLS all & Post-DS 2SLS & DS-fit IV \\
\midrule
Crime is a top-two national problem & 36.1 & 4.08 & 18.61 & 8.18 & 8.84 & 9.65 \\
 & & (1.68) & (6.79) & (3.37) & (3.88) & (4.36) \\
 & & [0.017] & [0.006] & [0.015] & [0.023] & [0.027] \\
\addlinespace
Crime is a top-two issue affecting me & 34.9 & 5.39 & 14.94 & 8.43 & 9.32 & 11.40 \\
 & & (2.00) & (7.00) & (4.31) & (4.73) & (5.05) \\
 & & [0.008] & [0.033] & [0.050] & [0.049] & [0.024] \\
\addlinespace
Crime affects my quality of life & 63.2 & 6.72 & 16.07 & 9.95 & 11.12 & 13.73 \\
 & & (2.09) & (6.67) & (3.92) & (4.25) & (4.69) \\
 & & [0.002] & [0.016] & [0.011] & [0.009] & [0.003] \\
\addlinespace
Feels unsafe & 17.4 & 6.27 & 4.79 & 6.81 & 6.28 & 6.59 \\
 & & (1.93) & (6.49) & (3.15) & (3.49) & (3.86) \\
 & & [0.002] & [0.460] & [0.031] & [0.072] & [0.088] \\
\addlinespace
Expects to be a victim & 43.8 & 6.59 & 16.91 & 15.44 & 15.03 & 16.98 \\
 & & (2.61) & (8.92) & (4.67) & (4.95) & (5.30) \\
 & & [0.013] & [0.058] & [0.001] & [0.002] & [0.001] \\
\addlinespace
Summary index (first principal component) & 39.4 & 5.84 & 13.58 & 9.40 & 9.64 & 10.97 \\
 & & (1.52) & (5.32) & (2.79) & (3.22) & (3.65) \\
 & & [0.000] & [0.011] & [0.001] & [0.003] & [0.003] \\
\bottomrule
\end{tabular}
\begin{tablenotes}[flushleft]
\footnotesize
\item Notes: $N = 101$ municipalities; 2008--2017 long differences. Coefficients are the
effect of doubling the immigrant share of the population, in percentage points. Controls:
intercept, municipality average age and share male in 2017. OLS: HC1 standard errors and
$t$ $p$-values. IV columns: HC0 standard errors and normal $p$-values (as in the paper's
\texttt{ivreg2}). Post-DS 2SLS standard errors treat the selected instruments as fixed.
Variables: \texttt{del\_problema3}, \texttt{del\_afecta3}, \texttt{del\_calvida3},
\texttt{del\_inseg5}, \texttt{del\_vict}, \texttt{pc\_19}.
\end{tablenotes}
\end{threeparttable}
\end{table}

\paragraph{Crime-related concerns.} The paper's central result survives: immigration
significantly raises crime-related concerns under every estimator. The effects are smaller,
though. The post-Dantzig estimates are roughly 50--70\% of the Bartik estimates; for
example, the summary index rises by 9.6 rather than 13.6 percentage points when the
immigrant share doubles (mean 39.4). The standard errors fall by about 30--45\%, so
precision improves even as the magnitudes shrink. ``Expects to be a victim,'' only
marginally significant in the paper, becomes clearly significant.

The estimates now sit between OLS and the Bartik 2SLS, which has two possible readings.
The Bartik estimates may have been noisy and somewhat inflated. Alternatively, the
overidentified estimators may be pulled toward OLS by many-instrument bias. The DS-fit IV
helps distinguish these. It is just-identified, so it is not subject to that bias, yet it
gives similar or slightly larger numbers (10.97 for the index). That suggests the gap from
the Bartik estimate is not purely a many-instrument artifact.

% ---------------------------------------------------------------- Table VII
\begin{table}[htbp]
\centering
\small
\begin{threeparttable}
\caption{Beliefs about crime trends (paper's Table VII)}
\label{tab:beliefs}
\begin{tabular}{lcccccc}
\toprule
 & & \multicolumn{2}{c}{Paper's design} & \multicolumn{3}{c}{Shares as separate instruments}\\
\cmidrule(lr){3-4}\cmidrule(lr){5-7}
Outcome & Mean & OLS & 2SLS Bartik & 2SLS all & Post-DS 2SLS & DS-fit IV \\
\midrule
Crime rising in my neighborhood & 42.1 & 9.75 & 13.44 & 8.90 & 8.70 & 9.64 \\
 & & (2.57) & (8.27) & (4.74) & (5.09) & (5.43) \\
 & & [0.000] & [0.104] & [0.060] & [0.087] & [0.076] \\
\addlinespace
Crime rising in my municipality & 64.9 & 6.94 & 8.74 & 6.70 & 6.62 & 5.63 \\
 & & (2.38) & (8.65) & (4.24) & (4.80) & (5.20) \\
 & & [0.004] & [0.312] & [0.114] & [0.168] & [0.278] \\
\addlinespace
Crime rising in the country & 78.9 & 2.73 & 1.71 & 2.25 & 0.78 & $-0.36$ \\
 & & (1.19) & (4.82) & (2.06) & (2.25) & (2.54) \\
 & & [0.023] & [0.722] & [0.274] & [0.730] & [0.886] \\
\bottomrule
\end{tabular}
\begin{tablenotes}[flushleft]
\footnotesize
\item Notes: See Table~\ref{tab:concerns}. Variables: \texttt{del\_tend\_barrio},
\texttt{del\_tend\_com}, \texttt{del\_tend\_pais}.
\end{tablenotes}
\end{threeparttable}
\end{table}

\paragraph{Beliefs about crime trends.} As in the paper, the IV estimates are positive for
local beliefs but not significant at conventional levels. The neighborhood-level estimate
is marginal ($p \approx 0.06$--$0.09$) under all three share-based estimators. Effects are
largest for the neighborhood and shrink toward zero for the country as a whole, consistent
with immigration shaping perceptions of one's \emph{local} environment.

% ---------------------------------------------------------------- Table VIII
\begin{table}[htbp]
\centering
\small
\begin{threeparttable}
\caption{Crime-preventive behavior (paper's Table VIII)}
\label{tab:prevent}
\begin{tabular}{lcccccc}
\toprule
 & & \multicolumn{2}{c}{Paper's design} & \multicolumn{3}{c}{Shares as separate instruments}\\
\cmidrule(lr){3-4}\cmidrule(lr){5-7}
Outcome & Mean & OLS & 2SLS Bartik & 2SLS all & Post-DS 2SLS & DS-fit IV \\
\midrule
Invests in home security (index) & 22.8 & 1.84 & 10.03 & 6.41 & 6.86 & 6.67 \\
 & & (1.09) & (4.70) & (2.32) & (2.61) & (2.78) \\
 & & [0.093] & [0.033] & [0.006] & [0.009] & [0.016] \\
\addlinespace
Coordinates with neighbors (index) & 13.2 & 0.80 & 12.44 & 5.72 & 6.77 & 8.09 \\
 & & (0.95) & (4.06) & (1.67) & (1.98) & (2.25) \\
 & & [0.399] & [0.002] & [0.001] & [0.001] & [0.000] \\
\addlinespace
Household owns a firearm & 4.8 & 0.45 & 0.92 & 0.34 & 0.65 & $-0.24$ \\
 & & (0.75) & (1.62) & (0.93) & (0.95) & (1.09) \\
 & & [0.548] & [0.571] & [0.714] & [0.496] & [0.825] \\
\addlinespace
Summary index (first principal component) & 16.4 & 1.38 & 11.44 & 6.13 & 6.90 & 7.41 \\
 & & (0.84) & (4.04) & (1.83) & (2.13) & (2.32) \\
 & & [0.104] & [0.005] & [0.001] & [0.001] & [0.001] \\
\bottomrule
\end{tabular}
\begin{tablenotes}[flushleft]
\footnotesize
\item Notes: See Table~\ref{tab:concerns}. Variables: \texttt{Index\_Vivienda},
\texttt{Index\_Vecinos}, \texttt{del\_arma}, \texttt{pc\_20}.
\end{tablenotes}
\end{threeparttable}
\end{table}

\paragraph{Crime-preventive behavior.} The pattern matches Table~\ref{tab:concerns}.
Immigration significantly increases investment in home security and coordination with
neighbors, but not gun ownership. The post-Dantzig estimates are about 55--70\% of the
Bartik estimates, with standard errors roughly half as large. Doubling the immigrant share
raises the preventive-behavior index by 6.9 percentage points (mean 16.4), compared with
11.4 in the paper.

% ---------------------------------------------------------------- Table IX
\begin{table}[htbp]
\centering
\small
\begin{threeparttable}
\caption{Self-reported victimization (paper's Table IX)}
\label{tab:victim}
\begin{tabular}{lcccccc}
\toprule
 & & \multicolumn{2}{c}{Paper's design} & \multicolumn{3}{c}{Shares as separate instruments}\\
\cmidrule(lr){3-4}\cmidrule(lr){5-7}
Outcome & Mean & OLS & 2SLS Bartik & 2SLS all & Post-DS 2SLS & DS-fit IV \\
\midrule
Larceny / theft & 8.5 & 0.93 & 0.32 & 1.87 & 1.46 & 0.94 \\
 & & (1.07) & (3.56) & (1.64) & (1.76) & (1.91) \\
 & & [0.384] & [0.928] & [0.255] & [0.407] & [0.621] \\
\addlinespace
Robbery by surprise (snatching) & 4.6 & $-0.03$ & $-0.60$ & $-0.89$ & $-0.60$ & $-0.68$ \\
 & & (0.73) & (2.95) & (1.15) & (1.35) & (1.39) \\
 & & [0.967] & [0.838] & [0.438] & [0.658] & [0.626] \\
\addlinespace
Vehicle theft & 0.8 & 0.34 & $-0.95$ & $-0.08$ & $-0.28$ & $-0.09$ \\
 & & (0.26) & (0.71) & (0.41) & (0.43) & (0.47) \\
 & & [0.196] & [0.181] & [0.846] & [0.513] & [0.852] \\
\addlinespace
Burglary & 4.8 & 2.03 & 1.13 & 2.64 & 2.79 & 3.14 \\
 & & (0.64) & (1.83) & (0.95) & (0.92) & (0.92) \\
 & & [0.002] & [0.538] & [0.005] & [0.002] & [0.001] \\
\addlinespace
Assault & 1.9 & 0.05 & 1.87 & 1.19 & 0.97 & 1.66 \\
 & & (0.60) & (1.59) & (0.79) & (0.89) & (0.93) \\
 & & [0.934] & [0.240] & [0.130] & [0.273] & [0.075] \\
\addlinespace
Violent robbery & 4.4 & 0.36 & 3.09 & 0.46 & 0.94 & 1.69 \\
 & & (0.66) & (2.10) & (0.89) & (0.98) & (1.11) \\
 & & [0.584] & [0.141] & [0.609] & [0.338] & [0.129] \\
\addlinespace
Total victimization & 21.5 & 3.03 & 3.50 & 4.62 & 4.73 & 5.70 \\
 & & (1.66) & (5.21) & (2.39) & (2.50) & (2.73) \\
 & & [0.072] & [0.501] & [0.053] & [0.059] & [0.037] \\
\bottomrule
\end{tabular}
\begin{tablenotes}[flushleft]
\footnotesize
\item Notes: See Table~\ref{tab:concerns}. Crime categories are our English labels for the
ENUSC categories. Variables: \texttt{vict\_hurto}, \texttt{vict\_robosorpr},
\texttt{vict\_robovehi}, \texttt{vict\_roboviv}, \texttt{vict\_lesiones},
\texttt{vict\_roboviolent}, \texttt{vict\_agreg}.
\end{tablenotes}
\end{threeparttable}
\end{table}

\paragraph{Victimization.} This is where our results depart most from the paper. Most
crime categories remain null, but \emph{burglary} becomes significant under all three
share-based estimators: doubling the immigrant share raises burglary victimization by about
2.8 percentage points (mean 4.8), with $p = 0.002$. That $p$-value would survive a Holm
correction across the table's seven outcomes. Total victimization is also marginally
positive ($p = 0.04$--$0.06$). The burglary effect is not new: it appears in the paper's
OLS and two-way fixed-effects results, where the authors treat it as not robust to
multiple testing, and the Bartik IV was too imprecise to detect it. It also does not depend
on selection, since the unselected all-shares 2SLS gives a similar estimate (2.64). The
crime-perception gap therefore remains: concerns rise by about 9.6 points against at most
about 4.7 points in total victimization. But it is narrower than in the paper (13.6 against
3.5).

% =====================================================================================
\section{Takeaways and next steps}

\paragraph{Takeaways.} The paper's qualitative conclusions largely hold up. Immigration
raises crime concerns and preventive behavior and has little effect on beliefs about
national crime trends. Using the shares as separate instruments, with Dantzig selection,
gives a first stage that explains about three times as much variation as the Bartik
instrument. The resulting estimates are smaller by roughly a third to a half and
noticeably more precise. They also suggest a modest effect on burglary that the paper's IV
could not detect.

\paragraph{Caveats.} The identifying variation now comes from the Argentine, Venezuelan,
and Chinese shares rather than the Peruvian share, so the paper's pre-trend evidence does
not directly validate our estimates. The post-Dantzig standard errors treat the selected
instruments as fixed. The choice of $\lambda$ assumes homoskedastic errors. Finally, we
have not yet applied the paper's multiple-testing corrections. Table X (homicides) is
omitted because its data are confidential and not in the replication package.

\paragraph{Next steps.}
\begin{itemize}
  \item Run the paper's pre-trend tests (Goldsmith-Pinkham, Sorkin and Swift 2020) for the
        Argentine, Venezuelan, and Chinese shares.
  \item Test overidentifying restrictions (Hansen $J$) and report weak-instrument-robust
        inference (Anderson--Rubin confidence sets, effective $F$).
  \item Compare penalization methods in the first stage: Lasso and post-Lasso with
        heteroskedasticity-robust penalty loadings (Belloni et al.\ 2012), the elastic net,
        and ridge. Choose $\lambda$ by cross-validation as well as by the theoretical rule.
  \item Add nonlinear first stages (random forests, boosting) with cross-fitting, and DML
        with flexible controls.
  \item Repeat the analysis with the work-permit and work-visa immigration measures, and
        apply Holm corrections within each table.
\end{itemize}

\end{document}